\section{Two exact very-high-prime Gram forms}
\label{sec:tpc45-affine-gram}

There are two natural operations on the signs below the cutoff \(y\):
one may set them to their physical value \(-1\), or one may average
them.  The first operation retains every cross-label inner product and
gives an affine polynomial in the high signs.  The second keeps only
equal residual-label fibers and gives a pure quadratic Gram synthesis.
These are different identities in different Hilbert spaces.

Let \(\cP_y\) be the finite set of primes \(p>y\) occurring as the high
part of an actual equality label.  Use the factorization
\begin{equation}
 u=sq,
 \qquad \Pplus(s)\le y,
 \qquad q=1\ \text{or}\ q\in\cP_y
 \label{eq:tpc45-u-sq}
\end{equation}
from \cref{thm:tpc45-unique-very-high-prime}.

\subsection{Physical low signs: an affine output-space form}

Define vectors in the original output space \(\Hh_{\rm out}\) by
\begin{align}
 (v_0)_i
 &:=\sum_{\substack{u\in\cU_i\\u_{>y}=1}}
     \mu(u)b_{i,u},
 \label{eq:tpc45-v0}\\
 (v_p)_i
 &:=\sum_{\substack{u\in\cU_i\\u_{>y}=p}}
     \mu(u/p)b_{i,u},
 \qquad p\in\cP_y.
 \label{eq:tpc45-vp}
\end{align}
By the singleton assertion in
\cref{thm:tpc45-unique-very-high-prime}, the sum in
\eqref{eq:tpc45-vp} has at most one term for each fixed \(i\).

\begin{theorem}[Physical-low affine normal form]
\label{thm:tpc45-physical-low-affine}
After setting every sign at \(q\le y\) to \(-1\), the complete equality
column is
\begin{equation}
 \boxed{
 S(-\one_{\le y},\varepsilon_{>y})
 =v_0+\sum_{p\in\cP_y}\varepsilon_pv_p}
 \quad\text{in }\Hh_{\rm out}.
 \label{eq:tpc45-affine-column}
\end{equation}
Consequently,
\begin{equation}
 \boxed{
 Z_0(-\one_{\le y},\varepsilon_{>y})
 =\left\|v_0+\sum_{p\in\cP_y}\varepsilon_pv_p\right\|^2.}
 \label{eq:tpc45-affine-energy}
\end{equation}
Its high-prime Walsh spectrum has degree at most two.  Explicitly, for
the centered polynomial \(F=Z_0-\cD_0\), the coefficients after the low
restriction are
\begin{align}
 c_\varnothing
 &=\|v_0\|^2+\sum_p\|v_p\|^2-\cD_0,
 \label{eq:tpc45-affine-c0}\\
 c_p&=2\operatorname{Re}\langle v_0,v_p\rangle,
 \label{eq:tpc45-affine-cp}\\
 c_{pq}&=2\operatorname{Re}\langle v_p,v_q\rangle
 \qquad(p<q),
 \label{eq:tpc45-affine-cpq}
\end{align}
and every coefficient of degree at least three is zero.  At the physical
high corner,
\begin{equation}
 \boxed{Z_0(-\one)=\left\|v_0-\sum_{p\in\cP_y}v_p\right\|^2.}
 \label{eq:tpc45-affine-physical-corner}
\end{equation}
\end{theorem}

\begin{proof}
If \(u_{>y}=1\), its character at the restricted low face is
\(\mu(u)\).  If \(u=sp\), the same character is
\(\mu(s)\varepsilon_p\).  Grouping the coordinate formulas
\cref{eq:tpc45-left-coordinate-Walsh,eq:tpc45-right-coordinate-Walsh}
gives \eqref{eq:tpc45-affine-column}; taking the output norm gives
\eqref{eq:tpc45-affine-energy}.  Expansion of that norm proves
\cref{eq:tpc45-affine-c0,eq:tpc45-affine-cp,eq:tpc45-affine-cpq}.
Finally put every remaining sign equal to \(-1\).
\end{proof}

The formula is literal in the row variables.  If the unique left source
row at \((j,p)\) is \(\alpha_p(j)\) and
\begin{equation}
 n_{\alpha_p(j),j}=p r_p(j),
 \label{eq:tpc45-left-pr}
\end{equation}
then
\begin{equation}
 (v_p)_{L,\gamma,j}
 =\mu(d_{\alpha_p(j)}r_p(j))
  B^L_{\alpha_p(j);\gamma,j},
 \label{eq:tpc45-vp-left-actual}
\end{equation}
with value zero when no such row exists.  Similarly, if the unique right
summation row is \(\gamma_p(j)\), then
\begin{equation}
 (v_p)_{R,\alpha,j}
 =\mu(d_{\gamma_p(j)}r_p(j))
  B^R_{\gamma_p(j);\alpha,j}.
 \label{eq:tpc45-vp-right-actual}
\end{equation}
Thus the residual physical sign is precisely \(\mu(dr)\).  The high
prime itself supplies the final minus sign in
\eqref{eq:tpc45-affine-physical-corner}.

For example, the left part of the quadratic coefficient is
\begin{align}
 c_{pq}^L
 =2\operatorname{Re}\sum_{\gamma,j}
 &\mu(d_{\alpha_p(j)}r_p(j))
  \mu(d_{\alpha_q(j)}r_q(j))
 \notag\\
 &\times B^L_{\alpha_p(j);\gamma,j}
  \overline{B^L_{\alpha_q(j);\gamma,j}},
 \label{eq:tpc45-cpq-actual-left}
\end{align}
where a summand is present only when
\begin{equation}
 p r_p(j)-m_{\alpha_p(j)}j=h,
 \qquad
 q r_q(j)-m_{\alpha_q(j)}j=h.
 \label{eq:tpc45-two-determinant-ladders}
\end{equation}
In particular,
\begin{equation}
 p r_p(j)-q r_q(j)
 =(m_{\alpha_p(j)}-m_{\alpha_q(j)})j.
 \label{eq:tpc45-determinant-difference}
\end{equation}
The right contribution is obtained from
\eqref{eq:tpc45-vp-right-actual}.  This is an actual affine determinant
sum with the source weights and joint masks retained.

The degree-two reduction also simplifies the TPC--44 coefficient
correlation.  Process the primes in \(\cP_y\) increasingly, after all
low signs have already been restricted.  Just before processing \(p\),
put
\begin{equation}
 w_p:=v_0-\sum_{\substack{q\in\cP_y\\q<p}}v_q.
 \label{eq:tpc45-wp}
\end{equation}
The surviving coefficients are
\begin{align}
 c_\varnothing^{(p)}
 &=\|w_p\|^2+\sum_{q\ge p}\|v_q\|^2-\cD_0,
 \label{eq:tpc45-current-c0}\\
 c_q^{(p)}&=2\operatorname{Re}\langle w_p,v_q\rangle
 \qquad(q\ge p),
 \label{eq:tpc45-current-cq}\\
 c_{qr}^{(p)}&=2\operatorname{Re}\langle v_q,v_r\rangle
 \qquad(p\le q<r).
 \label{eq:tpc45-current-cqr}
\end{align}
Let \(C_{p^-}(m)\) denote the Walsh coefficient of \(\chi_m\) in this
current centered polynomial after all lower signs and all
\(q\in\cP_y\) with \(q<p\) have been fixed to \(-1\); the index
\(m=1\) denotes the constant coefficient
\(C_{p^-}(1)=c_\varnothing^{(p)}\).
Hence the general sum over residual kernels is exactly
\begin{equation}
 \boxed{
 \sum_m C_{p^-}(m)\overline{C_{p^-}(pm)}
 =c_\varnothing^{(p)}c_p^{(p)}
  +\sum_{q>p}c_q^{(p)}c_{pq}.}
 \label{eq:tpc45-top-paired-collapse}
\end{equation}
Indeed, a residual \(m\) with at least two prime factors would make
\(pm\) have degree at least three, whose coefficient vanishes.

\subsection{Averaged low signs: the residual-label Gram form}

Now use a different Hilbert space,
\begin{equation}
 \Hh_{\rm res}
 :=\ell^2\bigl\{(i,s):
 i\in\cI_L\sqcup\cI_R,
 \ s\text{ squarefree},\ \Pplus(s)\le y\bigr\}.
 \label{eq:tpc45-residual-Hilbert-space}
\end{equation}
For \(p\in\cP_y\), define
\begin{equation}
 (x_p)_{i,s}:=
 \begin{cases}
  b_{i,sp},&sp\in\cU_i,\\
  0,&sp\notin\cU_i.
 \end{cases}
 \label{eq:tpc45-xp}
\end{equation}
Let
\begin{equation}
 \cD_H:=\sum_{p\in\cP_y}\sum_i
        \sum_{\substack{u\in\cU_i\\u_{>y}=p}}|b_{i,u}|^2.
 \label{eq:tpc45-DH}
\end{equation}

\begin{theorem}[Low-averaged pure Gram identity]
\label{thm:tpc45-low-averaged-Gram}
Assume both cutoff conditions
\eqref{eq:tpc45-y-geometric-conditions} and
\eqref{eq:tpc45-y-label-separation}.  Then
\begin{equation}
 \boxed{
 \E_{\varepsilon_q:q\le y}
 Z_0(\varepsilon_{\le y},-\one_{>y})
 =\cD_0-\cD_H+
  \left\|\sum_{p\in\cP_y}x_p\right\|_{\Hh_{\rm res}}^2.}
 \label{eq:tpc45-low-averaged-Gram}
\end{equation}
Moreover,
\begin{equation}
 \sum_p\|x_p\|^2=\cD_H,
 \qquad
 \|x_p\|^2=\|v_p\|^2,
 \label{eq:tpc45-xp-diagonal}
\end{equation}
and therefore
\begin{equation}
 \E_{\le y}Z_0(\varepsilon_{\le y},-\one_{>y})
 =\cD_0+2\operatorname{Re}
  \sum_{p<q}\langle x_p,x_q\rangle.
 \label{eq:tpc45-pure-quadratic-Gram}
\end{equation}
There is no linear term in \eqref{eq:tpc45-pure-quadratic-Gram}.
\end{theorem}

\begin{proof}
Fix an output coordinate \(i\).  A label with no prime above \(y\) is
its own low Walsh label \(u\), and distinct such full labels are
orthogonal under the low-sign average.  A label \(sp\) with \(p>y\)
has low Walsh label \(s\), while its fixed high character contributes
the sign \(-1\).

No no-high full label can have the same low label as a high-prime
residual.  Indeed, by \eqref{eq:tpc45-support-endpoints}, a no-high label
has size at least \(D_-N_-\), whereas
\begin{equation}
 s=\frac{d_\rho n_{\rho,j}}p
 \le\frac{D_+N_+}{y}<D_-N_-
 \label{eq:tpc45-residual-size-separation}
\end{equation}
for a high-prime label.  Thus low Walsh orthogonality leaves the atomic
energy \(\cD_0-\cD_H\) from the no-high labels and, for every low label
\(s\), the coefficient
\begin{equation}
 -\sum_{p\in\cP_y}b_{i,sp}
 \label{eq:tpc45-matched-high-coefficient}
\end{equation}
from the high-prime labels.  Squaring
\eqref{eq:tpc45-matched-high-coefficient} and summing \((i,s)\) gives
\eqref{eq:tpc45-low-averaged-Gram}.  Unique high-prime factorization
partitions the high atoms, proving the first identity in
\eqref{eq:tpc45-xp-diagonal}.  Fixed-\((j,p)\) singleton incidence shows
that \eqref{eq:tpc45-vp} contains at most one term per output coordinate;
multiplication by \(\mu(s)\) preserves its modulus, proving the second.
Expanding the last norm now gives
\eqref{eq:tpc45-pure-quadratic-Gram}.
\end{proof}

The M\"obius factors suppressed in the stripped definition of \(x_p\)
may be restored without changing the synthesis norm.  Indeed, the
diagonal map
\begin{equation}
 (Uy)_{i,s}:=\mu(s)y_{i,s}
 \label{eq:tpc45-residual-unitary-reweighting}
\end{equation}
is unitary on \(\Hh_{\rm res}\), and on every occurring factorization
\(s=dr\) one has \(\mu(s)=\mu(d)\mu(r)\).  Consequently
\begin{equation}
 \left\|\sum_p x_p\right\|
 =\left\|\sum_p Ux_p\right\|.
 \label{eq:tpc45-residual-physical-reweighting}
\end{equation}
This is the precise bridge from the pure Gram form to the source- and
cofactor-M\"obius coefficient slots discussed below; it does not impose
independence or cancellation on either factor.

The matching condition in this Gram space is completely explicit:
\begin{equation}
 \langle x_p,x_q\rangle
 =\sum_i\sum_{\substack{s\\sp,sq\in\cU_i}}
 b_{i,sp}\overline{b_{i,sq}}.
 \label{eq:tpc45-xp-xq-matching}
\end{equation}
For a left coordinate, a term in
\eqref{eq:tpc45-xp-xq-matching} consists of rows \(\alpha,\beta\) and
cofactors \(r_\alpha,r_\beta\) satisfying
\begin{align}
 p r_\alpha&=\ell_\alpha d_\alpha j+h,
 &q r_\beta&=\ell_\beta d_\beta j+h,
 \label{eq:tpc45-matching-ladders}\\
 d_\alpha r_\alpha&=d_\beta r_\beta=s,
 \label{eq:tpc45-residual-match}
\end{align}
and its weight is
\begin{equation}
 B^L_{\alpha;\gamma,j}
 \overline{B^L_{\beta;\gamma,j}}.
 \label{eq:tpc45-matching-weight-left}
\end{equation}
The right formula exchanges the two row systems exactly as in
\eqref{eq:tpc45-C-right}.

Because every active \(d\) and \(r\) in
\eqref{eq:tpc45-residual-match} is squarefree and coprime, one has
\begin{equation}
 d_\alpha\mid s,
 \qquad r_\alpha=s/d_\alpha,
 \qquad
 \ell_\alpha=\frac{ps-hd_\alpha}{d_\alpha^2j},
 \label{eq:tpc45-matching-row-recovery}
\end{equation}
and the same formula with \((q,d_\beta)\).  Thus a fixed
\((i,s,p)\)-fiber has at most \(\tau(s)\) candidate source divisors;
integrality, primality, dyadic support, and the actual mask only delete
candidates.  This is the exact matching mechanism behind the residual
Gram space.

\begin{remark}[The two Gram matrices must not be identified]
\label{rem:tpc45-two-Grams-distinct}
The physical-low coefficient
\(\langle v_p,v_q\rangle_{\Hh_{\rm out}}\) sums every pair of
\(p\)- and \(q\)-rows sharing an output coordinate and weights it by
\(\mu(s_p)\mu(s_q)\).  The low-averaged coefficient
\(\langle x_p,x_q\rangle_{\Hh_{\rm res}}\) retains only the exact matches
\(s_p=s_q\); on such a match the two low M\"obius signs multiply to one.
Although \(\|v_p\|=\|x_p\|\), their cross inner products are different.
The affine form \eqref{eq:tpc45-affine-energy} and the pure Gram form
\eqref{eq:tpc45-low-averaged-Gram} therefore serve different descent
routes.
\end{remark}

The remaining estimate in either route is coefficient-specific.  In
the physical-low form it is an angle bound for the actual vectors
\(v_p\), hence for the signed determinant sums
\eqref{eq:tpc45-cpq-actual-left}.  In the low-averaged form it is a
Bessel bound for the matching synthesis
\begin{equation}
 \left\|\sum_{p\in\cP_y}x_p\right\|^2
 \label{eq:tpc45-residual-synthesis-gate}
\end{equation}
on the ladders \eqref{eq:tpc45-matching-ladders}--
\eqref{eq:tpc45-residual-match}.  Singleton prime incidence and the
diagonal identity \eqref{eq:tpc45-xp-diagonal} do not by themselves
control either coherent sum.
